\SourcePage{239}{242}
\begin{center}\textbf{P r o b l e m s}\end{center}

\ProblemHeading{1.} In the quasiclassical approximation, to exponential
accuracy, determine the probability of deuteron breakup in a collision with a
heavy nucleus treated as a fixed centre of a Coulomb field (E.~M. Lifshitz,
1939).

\SolutionHeading{Solution.} Collisions with zero orbital angular momentum
make the largest contribution to the reaction probability. In the
quasiclassical approximation these are ``head-on'' collisions, in which the
motion of the particles reduces to one-dimensional motion.

Let $E$ be the deuteron energy measured in units of $\varepsilon$, its
proton--neutron binding energy, and let $E_n$ and $E_p$ be the energies of the
released neutron and proton (in the same units). Introduce also the
dimensionless coordinate $q=r/(Ze^2/\varepsilon)$ (where $Ze$ is the nuclear
charge), and denote its generally complex value at the ``transition point,''
that is, at the ``instant of breakup'' of the deuteron, by $q_0$. Write
$E_n$, $E_p$, and $E$ in the form
\begin{equation}
 E_n=v_n^2/2,\qquad E_p=v_p^2/2+1/q_0,\qquad E=v_d^2+1/q_0;
 \tag{1}
\end{equation}
$v_n$, $v_p$, and $v_d$ are the particle ``velocities'' at the instant of
breakup, measured in units of $\sqrt{\varepsilon/m}$ (where $m$ is the nucleon
mass). Here $v_n$ is real and equals the velocity of the released neutron,
whereas $v_p$ and $v_d$ are complex. Conservation of energy and momentum at
the transition point gives
\begin{equation}
 E_p+E_n=E-1,\qquad v_p+v_n=2v_d,
 \tag{2}
\end{equation}
whence
\[
 v_p=2i+v_n,\qquad v_d=i+v_n,\qquad
 1/q_0=E+1-v_n^2+2iv_n.
\]
Before the transition, the action of the system corresponds to motion of the
deuteron in the nuclear field up to the breakup point; its imaginary part is
\begin{equation}
 \begin{aligned}
 \operatorname{Im}S_1
 &=Ze^2\sqrt{\frac m\varepsilon}\operatorname{Im}
   \int_\infty^{q_0}\sqrt{4\left(E-\frac1q\right)}\,dq={}\\
 &=Ze^2\sqrt{\frac m\varepsilon}\operatorname{Im}
 \left\{2q_0v_d-\frac2{\sqrt E}\cosh^{-1}\sqrt{q_0E}\right\}.
 \end{aligned}
 \tag{3}
\end{equation}
After the transition, the action corresponds to motion of the neutron and
proton away from the breakup point:

\SourcePage{240}{243}
\begin{equation}
 \begin{aligned}
 \operatorname{Im}S_2
 &=Ze^2\sqrt{\frac m\varepsilon}\operatorname{Im}
 \left\{\int_{q_0}^\infty v_n\,dq+
 \int_{q_0}^\infty\sqrt{2\left(E_p-\frac1q\right)}\,dq\right\}={}\\
 &=Ze^2\sqrt{\frac m\varepsilon}\operatorname{Im}
 \left\{-v_nq_0-v_pq_0+\sqrt{\frac2{E_p}}
 \cosh^{-1}\sqrt{q_0E_p}\right\}.
 \end{aligned}
 \tag{4}
\end{equation}
According to (52.1), the probability of the process is
\begin{equation}
 w\sim\exp\left\{-\frac{2Ze^2}\hbar\sqrt{\frac m\varepsilon}
 \operatorname{Im}\left[
 \sqrt{\frac2{E_p}}\cosh^{-1}\sqrt{q_0E_p}
 -\frac2{\sqrt E}\cosh^{-1}\sqrt{q_0E}\right]\right\}.
 \tag{5}
\end{equation}
In accordance with the origin of the first and second $\cosh^{-1}$
terms in the square brackets from expressions (4) and (3), the signs of their
imaginary parts must agree with the signs of $\operatorname{Im}v_p$ and
$\operatorname{Im}v_d$, respectively (the signs of the latter quantities in
solving equations (2) were chosen so that
$\operatorname{Im}(S_1+S_2)>0$ results).

Because $w$ depends exponentially on $E_n$, the total breakup probability
(including all values of $E_n$ and $E_p=E-1-E_n$) is determined by the
smallest absolute value of the exponent as a function of $E_n$. Analysis shows
that this value is attained as $E_n\to0$. Then $q_0=1/(E+1)$, and (5) gives
\[
 w\sim\exp\left\{-\frac{2Ze^2}\hbar\sqrt{\frac m\varepsilon}
 \left[\sqrt{\frac2{E-1}}\arccos\sqrt{\frac{E-1}{E+1}}
 -\frac2{\sqrt E}\arccos\sqrt{\frac E{E+1}}\right]\right\}.
\]
The condition for this formula to apply is that the exponent be large compared
with unity.

By calculating the imaginary part of the action $S=S_1+S_2$ at nonzero values
of $E_n$, one can find the energy distribution of the released particles. Near
$E_n=0$ we have\footnote{At $E_n=0$ the function
$\operatorname{Im}S(E_n)$ has a cusp, from which it increases in both
directions---for positive and negative $E_n$ (values $E_n<0$ would correspond
to capture of the neutron by the nucleus).}
\[
 \operatorname{Im}S(E_n)-\operatorname{Im}S(0)\approx
 E_n\left(\frac{d\operatorname{Im}S}{dE_n}\right)_{E_n=0}.
\]
Evaluation of the derivative gives
\[
 \frac{dw}{dE_n}\sim\exp\left\{-\frac{2Ze^2}\hbar
 \sqrt{\frac m\varepsilon}E_n\left[
 \frac{3-E}{(E-1)(E+1)^2}+
 \frac1{\sqrt{2(E-1)^3}}\arccos\sqrt{\frac{E-1}{E+1}}
 \right]\right\}.
\]

\ProblemHeading{2.} Determine the coefficient of over-barrier reflection at
particle energies for which perturbation theory is applicable.

\SolutionHeading{Solution.} The result follows from formula (43.1), in which
the initial and final wave functions are plane waves propagating in opposite
directions and normalized, respectively, to unit flux density and to a delta
function of momentum divided by $2\pi\hbar$. Here $d\nu=dp'/2\pi\hbar$, where
$p'$ is the momentum after reflection. Carrying out the integration over
$dp'$ in (43.1), with the delta function taken into account, gives

\SourcePage{241}{244}
\begin{equation}
 R=\frac{m^2}{\hbar^2p^2}
 \left|\int_{-\infty}^{\infty}U(x)e^{2ipx/\hbar}\,dx\right|^2.
 \tag{1}
\end{equation}
This formula is valid when the condition for applicability of perturbation
theory is satisfied: $Ua/\hbar v\ll1$, where $a$ is the width of the barrier
(see the note on p.~205), and at the same time $pa/\hbar\lesssim1$. The latter
condition ensures that the dependence $R(p)$ is not exponential; otherwise,
the applicability of formula (1) requires further investigation.

\ProblemHeading{3.} Determine the coefficient of over-barrier reflection from
a quasiclassical barrier when the function $U(x)$ has a kink at $x=x_0$.

\SolutionHeading{Solution.} If $U(x)$ has a singular feature at a real value
of $x$, the reflection coefficient is determined mainly by the field near
that point. To calculate it, perturbation theory may formally be applied
without requiring its applicability condition to hold for every $x$; it is
enough that the quasiclassical condition be satisfied. We thus arrive at
formula (1) of Problem~2, with the sole difference that the momentum of the
incident particle in it must be replaced by the value of $p(x)$ at $x=x_0$.

Choosing the kink as the point $x=0$, we have near it
\[
 U=-F_1x\quad\text{for }x>0,\qquad
 U=-F_2x\quad\text{for }x<0,
\]
with different values of $F_1$ and $F_2$. The integration over $dx$ is
performed by introducing the damping factor $e^{\pm\lambda x}$ into the
integrand (after which $\lambda\to0$ is taken). The result is
\[
 R=\frac{m^2\hbar^2}{16p_0^6}(F_2-F_1)^2,
\]
where $p_0=p(0)$.
